\section{Experimental Setup}
\label{sec:setup}

\subsection{Backbones and Inference Configuration}
\label{sec:llm}

All agents in a run share one backbone LLM. To test whether the findings depend on the backbone we use four OpenAI models that differ in capability: \texttt{gpt-3.5-turbo} (snapshot \texttt{gpt-3.5-turbo-0125}), \texttt{gpt-4o-mini} (snapshot \texttt{gpt-4o-mini-2024-07-18}), \texttt{gpt-4.1-nano} (snapshot \texttt{gpt-4.1-nano-2025-04-14}), and \texttt{gpt-4.1-mini} (snapshot \texttt{gpt-4.1-mini-2025-04-14}).
All models are accessed through the OpenAI Chat Completions API.
Every call uses temperature \(1.0\), the API defaults for nucleus sampling and output length, a fresh random seed, and \texttt{logprobs} enabled; the token log-probabilities are used only to compute the reviewer reliability score of Section~\ref{sec:wm}.
No system prompt, tool, or retrieval component is used.
All coordination strategies and all baselines use the same backbone, the same decoding configuration, and the same answer format within a comparison, so that differences in accuracy are attributable to the coordination strategy.
The reviewer-style configurations use one author, \(k=2\) reviewers, one meta-reviewer, and at most \(N=3\) review rounds.

\subsection{Benchmarks}

We evaluate on eight reasoning benchmarks (Table~\ref{tab:dataset_overview}):
\textbf{GSM8K}~\cite{cobbe2021gsm8k}, grade-school math word problems;
\textbf{SVAMP}~\cite{patel2021svamp}, arithmetic word problems with structural perturbations;
\textbf{GSM-Hard}~\cite{gao2023pal}, GSM8K problems with numbers replaced by large values;
\textbf{CIAR}~\cite{liang2024mad}, counter-intuitive arithmetic reasoning problems;
\textbf{StrategyQA}~\cite{geva2021strategyqa}, yes/no questions that require implicit multi-hop reasoning;
and three subsets of \textbf{MMLU}~\cite{hendrycks2021mmlu}: \emph{elementary mathematics}, \emph{moral scenarios}, and \emph{professional law}.
The first four benchmarks have numeric answers (scored correct within an absolute tolerance of \(0.01\)); the others are multiple-choice or binary and are scored by exact match of the selected option.

For each benchmark we fix one list of 50 question indices (all 48 questions for CIAR), sampled once, and evaluate \emph{every} method and \emph{every} backbone on exactly these questions, giving 398 questions per backbone.
The three reviewer-style configurations are run with all four backbones in two independent runs (1{,}592 paired questions, 3{,}184 runs per configuration in total), the external baselines once with \texttt{gpt-4.1-mini} and \texttt{gpt-4.1-nano}, and the ablations once with \texttt{gpt-4o-mini}, the backbone on which the centralized gate had the largest effect in the first run.
A disjoint set of 30 further questions per benchmark (210 in total; CIAR has no unused questions) is reserved as a development split for the exploratory variants of Sections~\ref{sec:pilot}--\ref{sec:combined}.

\begin{table}[!t]
\caption{Evaluation Benchmarks}
\label{tab:dataset_overview}
\centering
\footnotesize
\setlength{\tabcolsep}{3pt}
\begin{tabular}{llll}
\toprule
Benchmark & Domain & Answer & \#Q \\
\midrule
GSM8K & Math word problems & Numeric & 50 \\
SVAMP & Math word problems & Numeric & 50 \\
GSM-Hard & Math, large numbers & Numeric & 50 \\
MMLU-Elem. & Elementary mathematics & 4-way MC & 50 \\
CIAR & Counter-intuitive arithmetic & Numeric & 48 \\
StrategyQA & Implicit multi-hop QA & Binary & 50 \\
MMLU-Moral & Moral scenarios & 4-way MC & 50 \\
MMLU-Law & Professional law & 4-way MC & 50 \\
\bottomrule
\end{tabular}
\end{table}

\subsection{Compared Methods}

\paragraph{Reviewer-style configurations.}
\emph{Reviewer pipeline} is the message-based author--reviewer--meta-reviewer pipeline of MARS~\cite{wang2025mars} without any world model; it accepts a candidate when the meta-reviewer decides ``right''.
\emph{Ours (Peer)} and \emph{Ours (Centralized)} add the world model of Section~\ref{sec:method} in the two coordination architectures.
The three configurations share all prompts.

\paragraph{External baselines.}
We re-implement six methods in the same evaluation pipeline, with the same backbone and answer format:
\emph{CoT}, a single chain-of-thought call~\cite{wei2022chain};
\emph{Self-Consistency} with five samples and majority voting~\cite{wang2023selfconsistency};
\emph{Self-Reflection}, one answer followed by one self-critique and revision~\cite{madaan2023selfrefine,shinn2023reflexion};
\emph{Multi-Agent Debate} with three agents and two rounds~\cite{du2024improving};
\emph{MAD} with an affirmative debater, a negative debater, and a judge that may stop the debate, for at most three rounds~\cite{liang2024mad};
and \emph{ReConcile} with three agents, at most three discussion rounds, and confidence-weighted voting~\cite{chen2024reconcile}.
ReConcile was proposed for a mixture of different LLMs; to keep the backbone fixed we use its single-backbone variant and omit its ``convincing samples'' component.
We do not include AgentVerse or GPTSwarm as baselines: the former recruits task-specific expert roles and the latter optimizes a graph on a training split, so neither can be reproduced faithfully in a setting with a fixed role set and no training data.

\paragraph{Ablations.}
Starting from the centralized configuration we remove one world-model component at a time:
(i) \emph{w/o shared state summary}: reviewers do not receive \(\sigma(s_t)\);
(ii) \emph{w/o transition validation}: \(V\) is skipped;
(iii) \emph{w/o trajectory risk}: \(R\) is not computed, so the risk is neither shown to the meta-reviewer nor used in the gate (the consensus and reliability conditions remain);
(iv) \emph{w/o WM-guided acceptance}: the gate of Section~\ref{sec:gate} is removed and acceptance follows the meta-reviewer, while state, validation, and risk reporting remain.

\subsection{Metrics and Statistical Analysis}

We report accuracy, the average number of tokens (prompt plus completion) per question, the average number of LLM calls per question, and the average wall-clock time per question.
Because individual benchmarks contain at most 50 questions, per-benchmark differences below roughly 15 percentage points cannot be resolved, and we base our conclusions on comparisons pooled over benchmarks. We accompany accuracy differences with paired tests: for two methods evaluated on the same questions we compute the per-question difference in correctness, a 95\% bootstrap confidence interval of its mean, and a two-sided paired sign-flip permutation test (20{,}000 permutations).
We do not correct for multiple comparisons and therefore treat per-dataset \(p\)-values as descriptive.
Every number in the tables and figures of this paper is generated by a single analysis script from the per-question run logs.
